% ch07b-notifications-chat
\section{الإشعارات والشات: إزاي بيشتغلوا فعلاً}

في الفصل السابق فكّكنا \en{SignalR}: ثلاثة \en{hubs} ترث قاعدةً واحدة مصادَقة
بـ\en{JWT}، ومجموعات مثل \code{user:<id>} و\code{conversation:<id>}، والدفع اللحظيّ
سطرٌ واحد عبر \code{IHubContext}. الآن نبني \emph{فوق} تلك الأساسات ميزتين
يستعملهما كلّ مستخدم كل يوم: \textbf{جرس الإشعارات} الذي يومض دون تحديث الصفحة،
و\textbf{الشات الخاصّ 1:1} بين طرفين. لن نُعيد شرح \en{SignalR} نفسه؛ سنتتبّع
بدلًا من ذلك رحلة \emph{الحدث} — من اللحظة التي يقع فيها شيء في النظام حتى تظهر
النتيجة على شاشة الطرف الآخر.

\begin{idea}[الفكرة: ميزتان حيّتان فوق طبقةٍ واحدة]
الجرس والشات يبدوان مختلفَين، لكنّهما يتشاركان الهيكل نفسه: \textbf{يحفظ النظام
الحقيقة أوّلًا في قاعدة البيانات}، ثمّ \textbf{يدفعها حيّةً} عبر \en{SignalR} لِمن
هو متّصل الآن. الحفظ يضمن ألّا يضيع شيء (المستخدم غير المتّصل يجده حين يعود)،
والدفع يضمن الفوريّة لِمن هو حاضر. كل ما يلي مجرّد تفصيل لهاتين الخطوتين.
\end{idea}

\subsection{نقطة الدخول: \code{DispatchAsync} خطوة بخطوة}

كل إشعار في النظام — تأكيد حجز، دفعة نجحت، رسالة جديدة — يمرّ عبر باب واحد اسمه
\code{INotificationDispatcher.DispatchAsync}. المُرسِل لا يعرف شيئًا عن اللغات أو
القنوات أو ساعات الهدوء؛ يقول فقط «أبلِغ هذا المستخدم بهذا النوع»، والموزّع يتكفّل
بالباقي بترتيب ثابت. هذا هو قلب الدالة (مقتصّ):

\begin{english}
\begin{lstlisting}[language={[Sharp]C}]
public async Task DispatchAsync(
    Guid recipientUserId, NotificationType type, NotificationParams parameters,
    string? deepLink, string? idempotencyKey, CancellationToken ct)
{
    // 1) Idempotency — Stripe/event replays must not double-insert.
    if (!string.IsNullOrEmpty(idempotencyKey)
        && await db.Notifications.AnyAsync(n => n.IdempotencyKey == idempotencyKey, ct))
        return;

    var content  = catalog.Render(type, parameters);                  // 2) EN + AR
    var silenced = await IsSilencedAsync(recipientUserId, type, ct);   // 3) DND?
    var channels = await ResolveChannelsAsync(recipientUserId, type, silenced, ct);
    // 4) build a Notification row per channel, SaveChanges, then DeliverAsync each.
}
\end{lstlisting}
\end{english}

هنشوف إيه سطرًا-بسطر:
\begin{itemize}
\item الخطوة \textbf{الأولى} حارس ضدّ التكرار: إن وُجد \code{idempotencyKey} وسبق أن
  خُزِّن إشعارٌ بنفس المفتاح، نخرج فورًا (\code{return}). هذا يمنع الإشعار المكرَّر
  حين تعيد \en{Stripe} إرسال الحدث نفسه مرّتين.
\item الخطوة \textbf{الثانية} \code{catalog.Render}: تحوّل النوعَ والمعاملات إلى نصّ
  فعليّ \textbf{بالعربية والإنجليزية معًا} (عنوان + نصّ لكل لغة). المُرسِل لا يكتب
  نصوصًا، بل يمرّر نوعًا وقيمًا فقط.
\item الخطوة \textbf{الثالثة} \code{IsSilencedAsync}: هل هذا المستخدم في وضع «عدم
  إزعاج» الآن؟ نتيجة \code{silenced} تسري على كل ما بعدها.
\item الخطوة \textbf{الرابعة} \code{ResolveChannelsAsync}: تقرّر القنوات (جرس داخل
  التطبيق؟ بريد؟)، ثمّ يُبنى صفّ \code{Notification} لكل قناة، يُحفَظ، ثمّ يُسلَّم.
\end{itemize}

\subsubsection{منع التكرار: الفحص المسبق وسباق الفهرس الفريد}

الحارس المسبق (\code{AnyAsync}) يكفي في الحالة العاديّة، لكنّه لا يكفي وحده لو وصل
حدثان بنفس المفتاح في اللحظة ذاتها: قد يمرّ كلاهما من الفحص قبل أن يحفظ أيٌّ منهما.
لذلك توجد شبكة أمان ثانية على مستوى قاعدة البيانات: فهرس فريد على
\code{(IdempotencyKey, Channel)}. من يصل ثانيًا يصطدم بالفهرس، فنترجم الاصطدام إلى
«لا عمل» (\en{no-op}) بدل خطأ:

\begin{english}
\begin{lstlisting}[language={[Sharp]C}]
try { await db.SaveChangesAsync(ct); }
catch (ConflictException) when (!string.IsNullOrEmpty(idempotencyKey))
{
    // A concurrent dispatch with the same key won the race; the unique
    // (IdempotencyKey, Channel) index rejected our rows — detach and no-op.
    db.Notifications.RemoveRange(rows);
    return;
}
\end{lstlisting}
\end{english}

هنشوف إيه: لو نجح حفظ الطرف الآخر أوّلًا، يرفض الفهرسُ صفوفَنا فتُرمى
\code{ConflictException}. نلتقطها \textbf{بشرط} وجود مفتاح تكرار (\code{when})، ثمّ
نفصل صفوفنا المرفوضة (\code{RemoveRange}) ونخرج بهدوء. النتيجة النهائيّة واحدة سواء
وصل الحدث مرّة أو ألف مرّة: إشعارٌ واحد فقط.

\begin{why}
لماذا حارسان لا حارس واحد؟ الفحص المسبق يوفّر عملًا في الحالة الشائعة (لا نحاول
الحفظ أصلًا)، لكنّه ليس ذرّيًّا. الفهرس الفريد ذرّيّ ومضمون من قاعدة البيانات لكنّه
يكلّف استثناءً. الجمع بينهما يعطي الأداء الجيّد \textbf{و} الصحّة تحت التزامن معًا —
وهو نمط الحماية بطبقتين نفسه الذي رأيناه في المدفوعات.
\end{why}

\subsection{«عدم الإزعاج»: مكتوم، ساعات هدوء، وسياسة \en{fail-open}}

الدالة \code{IsSilencedAsync} تجيب عن سؤال واحد: هل نكتم هذا المستخدم الآن؟ الجواب
«نعم» في حالتين: إمّا أنّه فعّل الكتم الكامل (\code{NotificationsMuted})، أو أنّ
اللحظة الحاليّة تقع داخل \textbf{ساعات الهدوء} (\code{QuietHours}) التي حدّدها.
وفيها ثلاث دقائق مهمّة:

\begin{itemize}
\item \textbf{بمنطقة المستخدم الزمنيّة}: نافذة الهدوء تُحسَب بتوقيت المستخدم لا
  بتوقيت الخادم، فـ«العاشرة مساءً» تعني مساءه هو.
\item \textbf{قد تلفّ حول منتصف الليل}: نافذة مثل \en{22:00 $\rightarrow$ 08:00}
  تبدأ في يوم وتنتهي في اليوم التالي.
\item \textbf{اختبار المستخدم يمرّ دائمًا}: النوع \code{SystemTest} يُعيد «غير مكتوم»
  فورًا — لأنّه الوسيلة التي يتحقّق بها المستخدم من إعداداته نفسها.
\end{itemize}

منطق «اللفّ حول منتصف الليل» هو النقطة الذكيّة، وهذا هو الكود:

\begin{english}
\begin{lstlisting}[language={[Sharp]C}]
private static bool InQuietWindow(TimeOnly start, TimeOnly end, string? tzId)
{
    var localNow = TimeOnly.FromDateTime(
        TimeZoneInfo.ConvertTime(DateTimeOffset.UtcNow, ResolveTimeZone(tzId)).DateTime);
    // The window may wrap past midnight (e.g. 22:00 -> 08:00).
    return start <= end
        ? localNow >= start && localNow < end    // same day
        : localNow >= start || localNow < end;   // 22:00 -> 08:00
}
\end{lstlisting}
\end{english}

هنشوف إيه:
\begin{itemize}
\item نحوّل «الآن» \en{UTC} إلى الوقت المحلّي للمستخدم عبر \code{ConvertTime} ثمّ
  نأخذ الجزء الزمنيّ فقط (\code{localNow}).
\item إن كانت \code{start <= end} فالنافذة داخل اليوم نفسه، فنطلب أن يقع «الآن»
  بينهما (\code{and}).
\item وإلّا فالنافذة ملتفّة حول منتصف الليل، فنطلب أن يقع «الآن» \textbf{بعد} البداية
  \textbf{أو قبل} النهاية (\code{or}) — فتشمل ليلة كاملة تعبر منتصف الليل.
\end{itemize}

\begin{secnote}
لاحظ فلسفة \code{IsSilencedAsync} كلّها: \textbf{تفشل مفتوحةً} (\en{fail-open}). لو
كان ملفّ المستخدم مفقودًا، أو مُعرِّف المنطقة الزمنيّة تالفًا، أو أيّ استثناء آخر —
تُرجِع «غير مكتوم» فيصل الإشعار عاديًّا. القاعدة: خطأٌ في منطق الكتم يجب ألّا
\textbf{يبتلع} إشعارًا مهمًّا. الأسوأ المقبول هو إزعاجٌ عابر، لا معلومة ضائعة.
\end{secnote}

\subsection{القنوات والدفع الحيّ للجرس}

بعد معرفة حالة الكتم، تقرّر \code{ResolveChannelsAsync} أين يذهب الإشعار. القاعدة
بسيطة: \code{InApp} (الجرس) \textbf{دائمًا}، و\code{Email} حسب تفضيلات المستخدم لكل
\emph{نوع} إشعار على حدة. أمّا تحت «عدم الإزعاج» فالاثنان يُكتَمان — \textbf{لكن}
صفّ \code{InApp} يظلّ يُخزَّن في قاعدة البيانات. الجدول يلخّص السلوك:

\begin{center}
\begin{tabular}{L{4.4cm} C{2.4cm} C{2.9cm} C{2.0cm}}
\toprule
\textbf{\color{spPrimary}الحالة} & \textbf{\color{spPrimary}صفّ \code{InApp}} & \textbf{\color{spPrimary}دفعة حيّة (\en{toast})} & \textbf{\color{spPrimary}\code{Email}} \\
\midrule
عاديّ، وتفضيل البريد مفعّل   & يُخزَّن & يُدفَع & يُرسَل \\
عاديّ، وتفضيل البريد معطَّل  & يُخزَّن & يُدفَع & لا \\
تحت «عدم الإزعاج»           & يُخزَّن & لا     & لا \\
\bottomrule
\end{tabular}
\end{center}

المفتاح في السطر الأخير: تحت الكتم تُرجِع \code{ResolveChannelsAsync} قناة
\code{InApp} وحدها (تُسقِط \code{Email})، وفي داخل التسليم تُكتَم أيضًا الدفعةُ الحيّة
لتلك القناة. النتيجة: المستخدم \emph{يجد} الإشعار في قائمته لاحقًا، لكنّه لا يُزعَج
به لحظتَها. وهذا هو سطر الدفع الحيّ داخل \code{DeliverAsync} لقناة \code{InApp}:

\begin{english}
\begin{lstlisting}[language={[Sharp]C}]
// The in-app row is always stored; the live toast is suppressed under DND.
if (!silenced)
    await hub.Clients.Group($"user:{row.RecipientUserId}")
        .SendAsync("notification", new {
            id = row.Id, type = row.Type.ToString(),
            row.TitleEn, row.TitleAr, row.BodyEn, row.BodyAr,
            row.DeepLink, createdAt = DateTimeOffset.UtcNow }, ct);
\end{lstlisting}
\end{english}

هنشوف إيه: الشرط \code{if (!silenced)} هو ما يكتم الدفعة الحيّة فقط (الصفّ خُزِّن قبل
هذا السطر أصلًا). نستهدف مجموعة \code{user:<id>} — التي انضمّ إليها المستخدم عند
الاتّصال — ونرسل حدثًا اسمه \code{notification} يحمل العنوان والنصّ باللغتين مع
رابط عميق (\code{DeepLink}). الواجهة تستمع لهذا الاسم فيومض الجرس فورًا.

\begin{qa}
\textbf{سؤال متوقّع:} «تحت عدم الإزعاج، لماذا نخزّن صفّ الجرس أصلًا بدل تجاهله؟» لأنّ
الكتم يعني «لا تُزعِجني الآن»، لا «احذف الخبر». المستخدم قد يفتح قائمة الإشعارات بعد
ساعات الهدوء ويريد أن يجد كل ما فاته. الكتم يقمع \emph{المقاطعة} لا \emph{المعلومة}.
\end{qa}

\begin{keyp}[نقطة مهمة: لماذا \code{await} في \code{OnConnectedAsync}؟]
رأينا في الفصل السابق أنّ \code{NotificationHub} يضيف الاتّصال إلى مجموعة
\code{user:<id>} عند الاتّصال، \textbf{وينتظر} الإضافة (\code{await}) لا يطلقها
ويمضي. السبب حرِج هنا تحديدًا: لو دُفِع إشعارٌ في اللحظة التالية للاتّصال مباشرةً،
ومجموعةُ المستخدم لم تكتمل بعد، لَسَبَق الإشعارُ الانضمامَ وضاع على صاحبه — ثمّ ظهر
لاحقًا كأنّه «مكرَّر» بعد إعادة الإرسال. الانتظار يقفل هذه الفجوة الزمنيّة.
\end{keyp}

\subsection{الشات 1:1: الحضور والبوّابة قبل أي رسالة}

ننتقل للميزة الثانية. \code{ChatHub} يرث \code{AuthenticatedHub} المعلَّم
بـ\code{[Authorize]}، فلا يتّصل به إلّا حاملُ \en{JWT} صالح. فوق ذلك يتتبّع
\textbf{الحضور} (\en{online/offline}) بعدّاد مرجعيّ: أوّل اتّصال للمستخدم يعلنه
«متّصلًا» للآخرين، وآخر اتّصال يسقط (وصل العدّاد صفرًا) يعلنه «غير متّصل». والعميلُ
حديثُ الاتّصال يستدعي \code{GetOnlineUsers} ليهيّئ حالته دفعةً واحدة بدل انتظار
الأحداث تتقاطر.

لكن قبل أن ينضمّ أي عميل إلى مجموعة محادثة أو يبثّ مؤشّر كتابة، تمرّ كل عمليّة عبر
بوّابة واحدة، \code{EnsureParticipantAsync}:

\begin{english}
\begin{lstlisting}[language={[Sharp]C}]
private async Task EnsureParticipantAsync(string conversationId)
{
    if (!Guid.TryParse(Context.UserIdentifier, out var userId))
        throw new HubException("Unauthorized.");
    if (!Guid.TryParse(conversationId, out var convId))
        throw new HubException("Invalid conversation id.");

    var isParticipant = await db.Conversations.AsNoTracking()
        .AnyAsync(c => c.Id == convId
            && (c.ParticipantOneId == userId || c.ParticipantTwoId == userId));
    if (!isParticipant)
        throw new HubException("You are not a participant in this conversation.");
}
\end{lstlisting}
\end{english}

هنشوف إيه: نأخذ هويّة المتّصل من \code{Context.UserIdentifier} (من الـ\en{JWT} لا من
طلب العميل)، ثمّ نسأل قاعدة البيانات: هل هذا المستخدم أحد طرفَي المحادثة
(\code{ParticipantOneId} أو \code{ParticipantTwoId})؟ إن لم يكن، نرمي
\code{HubException} فورًا فلا انضمام ولا كتابة. الدوالّ \code{JoinConversation}
و\code{TypingStart} و\code{TypingStop} كلّها تستدعي هذه البوّابة أوّلًا.

\begin{secnote}
هذه بوّابة \en{SEC-02 IDOR}. بدونها كان \textbf{أي} مستخدم مصادَق يستطيع الانضمام
إلى \code{conversation:<id>} لأي محادثة — بمجرّد تخمين المُعرِّف — والتنصّت على
رسائل خاصّة تُدفَع لتلك المجموعة. الدرس: المصادقة (\en{authentication}) تُثبِت
\emph{مَن أنت}، لكنّ التفويض (\en{authorization}) — أنت طرفٌ في \emph{هذه}
المحادثة تحديدًا — يجب فحصه على الخادم قبل كل اشتراك، لا الاكتفاء بأنّك مسجَّل
دخولك.
\end{secnote}

مؤشّرات الكتابة تتبع النمط نفسه: \code{TypingStart}/\code{TypingStop} تبثّان — بعد
البوّابة — إلى \code{OthersInGroup} أي «كل مَن في المجموعة عداي». فيرى الطرف الآخر
«يكتب الآن…» دون أن يرى المرسِلُ إشعارَ نفسه.

\subsection{من الرسالة إلى الطرف الآخر: \code{SendMessageCommandHandler}}

إرسال رسالة يمرّ بستّ خطوات مرتّبة داخل معالِج واحد. المقتطف التالي يجمع أهمّها
(مقتصّ):

\begin{english}
\begin{lstlisting}[language={[Sharp]C}]
// 1) canonical order (smaller GUID first) -> one conversation per pair
var (participantOne, participantTwo) = senderId.CompareTo(request.RecipientId) <= 0
    ? (senderId, request.RecipientId)
    : (request.RecipientId, senderId);
// ... find-or-create the conversation with that ordered pair ...

// 2) block check (either direction)
if (isBlocked) throw new ConflictException("Cannot send message. User is blocked.");

// 3) SEC-07 defense-in-depth: strip HTML before it hits any sink
var sanitizedBody = new Ganss.Xss.HtmlSanitizer().Sanitize(request.Body);

// 4) persist the message + touch the conversation
db.Messages.Add(new ChatMessage { ConversationId = conversation.Id,
    SenderId = senderId, Body = sanitizedBody, SentAt = DateTimeOffset.UtcNow });
await db.SaveChangesAsync(ct);

// 5) live push to whoever has the conversation open
await chatNotifier.NotifyNewMessageAsync(
    conversation.Id, message.Id, senderId, sanitizedBody, message.SentAt, ct);
// 6) bell + email for the absent party (elsewhere in the app)
await publisher.Publish(new ChatMessageReceivedEvent(/* ... */), ct);
\end{lstlisting}
\end{english}

هنشوف إيه خطوةً-خطوة:
\begin{itemize}
\item \textbf{(1) الترتيب \en{canonical}}: نرتّب الطرفين بحيث يكون الأصغر
  (\code{CompareTo}) أوّلًا قبل إنشاء المحادثة. هذا يجعل الفهرس الفريد على الزوج
  يضمن \textbf{محادثة واحدة لكل ثنائيّ}، ويمنع سباق رسالتين متزامنتين من الطرفين
  يُنشئان صفّين بأعمدة مقلوبة.
\item \textbf{(2) فحص الحظر}: إن كان أيٌّ من الطرفين حاظرًا للآخر (في أيّ اتّجاه)
  نرمي \code{ConflictException} فلا رسالة تُرسَل.
\item \textbf{(3) تنقية \en{HTML}} عبر \code{HtmlSanitizer}: نجرّد أي وسوم أو سكربتات
  من النصّ \emph{قبل} تخزينه.
\item \textbf{(4) الحفظ}: نضيف \code{ChatMessage} ونحدّث المحادثة (آخر رسالة
  ووقتها)، ثمّ \code{SaveChanges}.
\item \textbf{(5) الدفع الحيّ}: \code{NotifyNewMessageAsync} يبثّ لمجموعة المحادثة.
\item \textbf{(6) الحدث}: ننشر \code{ChatMessageReceivedEvent} ليتكفّل بجرس الطرف
  الغائب وبريده.
\end{itemize}

\begin{secnote}
لماذا ننقّي \en{HTML} والواجهةُ \en{React} تعرض النصّ آمنًا أصلًا؟ لأنّ نفس النصّ
المخزَّن يذهب إلى \textbf{ثلاثة مصارف} (\en{sinks}): عقدة نصّ في \en{React}
(آمنة)، ودفعة \en{SignalR} حيّة، ومعاينة داخل بريد \en{HTML} — والأخيران ليسا
\en{React}. التنقية \textbf{عند التخزين} تجعل كل مصرف آمنًا مرّة واحدة بدل الاعتماد
على أن يتذكّر كلُّ مستهلِك أن يهرب النصّ. هذا معنى \en{SEC-07}: دفاعٌ في العمق.
\end{secnote}

الدفع الحيّ نفسه سطرٌ واحد في \code{ChatRealtimeNotifier}، يبثّ لمجموعة المحادثة:

\begin{english}
\begin{lstlisting}[language={[Sharp]C}]
hub.Clients.Group($"conversation:{conversationId}")
   .SendAsync("MessageReceived",
       new { conversationId, messageId, senderId, body, sentAt }, ct);
\end{lstlisting}
\end{english}

هنشوف إيه: نستهدف مجموعة \code{conversation:<id>} — التي لا ينضمّ إليها إلّا من
اجتاز بوّابة \code{EnsureParticipantAsync} — ونرسل حدثًا اسمه \code{MessageReceived}.
كلّ عميل يشاهد هذه المحادثة الآن يستقبله فتظهر الرسالة فورًا في الخيط.

\begin{keyp}[نقطة مهمة: المسار المزدوج]
لماذا خطوتان (5) و(6) معًا؟ لأنّهما تصلان إلى شخصين مختلفين. الدفع الحيّ
(\code{NotifyNewMessageAsync}) يصل \textbf{فقط} لِمن يفتح المحادثة على شاشته الآن —
هو مشترِكٌ في مجموعة \code{conversation:<id>}. أمّا الطرف الموجود في مكان آخر من
التطبيق (صفحة المنح مثلًا) فليس في تلك المجموعة، فيتكفّل به الحدثُ
\code{ChatMessageReceivedEvent}: يُطلِق إشعار جرسٍ عبر المسار الذي شرحناه في أوّل
الفصل. \textbf{الخلاصة}: المفتوحة أمامك تصلك حيّةً، والمغلقة يذكّرك بها الجرس — فلا
رسالة تفوت.
\end{keyp}

\begin{note}
لاحظ التناظر بين الميزتين: كلتاهما \textbf{تحفظ ثمّ تدفع}. الإشعار يحفظ صفّ
\code{Notification} ثمّ يدفع لمجموعة \code{user:<id>}؛ والرسالة تحفظ \code{ChatMessage}
ثمّ تدفع لمجموعة \code{conversation:<id>} \emph{و}تنشر حدثًا للجرس. الدوام أوّلًا،
الفوريّة ثانيًا — هو النمط الذي يجعل الميزتين موثوقتين لا مجرّد «حيّتين».
\end{note}

\begin{summary}[الخلاصة في سطور]
\begin{itemize}
\item كل إشعار يمرّ عبر \code{DispatchAsync} بترتيب ثابت: منع تكرار
  \en{$\rightarrow$} تصيير النصّ باللغتين (\code{catalog.Render}) \en{$\rightarrow$}
  فحص الكتم \en{$\rightarrow$} تحديد القنوات \en{$\rightarrow$} حفظ \en{$\rightarrow$}
  تسليم.
\item منع التكرار بطبقتين: فحص \code{IdempotencyKey} المسبق + فهرس فريد على
  \code{(IdempotencyKey, Channel)} يترجم السباق إلى \en{no-op}.
\item «عدم الإزعاج» = مكتوم أو داخل ساعات الهدوء (بمنطقة المستخدم، وقد تلفّ حول منتصف
  الليل)؛ السياسة \en{fail-open}، و\code{SystemTest} يمرّ دائمًا. تحت الكتم يُخزَّن
  صفّ \code{InApp} لكن لا \en{toast} ولا \code{Email}.
\item الشات محميّ ببوّابة \code{EnsureParticipantAsync} (\en{SEC-02 IDOR}) قبل أي
  انضمام أو كتابة — المتّصل من الـ\en{JWT} يجب أن يكون طرفًا في المحادثة.
\item \code{SendMessageCommandHandler}: ترتيب \en{canonical} (محادثة واحدة لكل زوج)
  \en{$\rightarrow$} فحص حظر \en{$\rightarrow$} تنقية \en{HTML} (\en{SEC-07}، لثلاثة
  مصارف) \en{$\rightarrow$} حفظ \en{$\rightarrow$} دفع حيّ \en{$\rightarrow$} حدثٌ
  للغائب.
\item المسار المزدوج: الدفع الحيّ لِمن يفتح المحادثة، وحدث الجرس لِمن هو في مكانٍ
  آخر — فلا رسالة تُفوَّت.
\end{itemize}
\end{summary}
